% ======================================================================
\section{Experiments: Selective Hamiltonian Risk Adaptation}
\label{sec:experiments}
% ======================================================================

The experiments ask one question across four settings: does Hamiltonian
context enrichment activate risk forces when a safer feasible alternative
exists and suppress them when it does not? We lead with the head-to-head
domain where this mechanism is decisive: delayed required escape, where
the policy must suppress a tempting lateral maneuver before it is feasible
and activate it after the escape opens, all with zero replans.

% ----------------------------------------------------------------------
\subsection{Setup, Metrics, and Domains}
\label{sec:exp_setup}
% ----------------------------------------------------------------------

\paragraph{Research questions.}
RQ1--RQ2 (spatial): does the context-enriched field activate toward a feasible lower-risk
direction and suppress when blocked?
RQ3--RQ5 (temporal): does the context-enriched field activate faster than replanning
baselines, avoid pursuing closed or unavailable escapes, and
suppress-then-activate within one episode?
RQ6: must risk enter the force channel, or is loss reweighting sufficient?
RQ7: does CVaR training reduce tail violations, not only mean risk?

\paragraph{Regimes.}
All domains share \textbf{R1} (feasible lower-risk alternative exists),
\textbf{R2} (lower-risk-looking direction blocked or dynamically unsafe),
and \textbf{R3} (risk-neutral; the context-enriched field should preserve geometry-only behavior).
RELLIS-Dyn extends R1/R2 from space to time via eight event types spanning
soft-risk, hard-boundary, dynamic-obstacle, and compound/delayed scenarios;
full definitions are in Appendix~\ref{app:rellis_dyn_details}. The primary
head-to-head benchmark is \emph{delayed required escape}: blocked before
$t_{\rm escape}$, necessary after.

\paragraph{Metrics.}
Spatial selectivity:
\begin{equation}
\mathrm{CAR}=\Pr[\langle F_{\rm ctx},d_{\rm safe}\rangle{>}\epsilon\mid R1],\;\;
\mathrm{FAR}=\Pr[|P_\perp F_{\rm ctx}|{>}\epsilon\mid R2],\;\;
\mathrm{SR}=\frac{\mathbb{E}[|P_\perp F_{\rm ctx}|\mid R1]}
                 {\mathbb{E}[|P_\perp F_{\rm ctx}|\mid R2]},
\end{equation}
and AUPRC (threshold-independent). Temporal: \emph{reaction delay}
($\delta$-displacement from the pre-event geometry-only trajectory; sensitivity in
Appendix Table~\ref{tab:app_rellis_dyn_delay_sensitivity}),
\emph{stale exposure} (violation before reaction),
\emph{false pre-activation} (lateral force before $t_{\rm escape}$),
\emph{suppress rate} ($1-$false pre-activation).
Post-event \emph{violation CVaR} (hard contact weighted separately) is
the primary tail metric.

% ----------------------------------------------------------------------
\subsection{Primary Head-to-Head: Delayed Required Escape}
\label{sec:exp_primary_delayed}
% ----------------------------------------------------------------------

All methods receive the same updated BEV patch at every step; only the
decision mechanism differs. Table~\ref{tab:delayed_escape} is the main
head-to-head result: the route-aware CVaR field reduces DWA and semantic
MPPI false pre-activation from $0.950$ to $0.180$ while raising success
from $0.480$/$0.240$ to $0.810$, and does so without replans. The table also isolates the
three ingredients: Hamiltonian force structure, CVaR timing, and
context-conditioned coefficients.

\begin{table}[t]
\centering
\caption{\textbf{Delayed-required-escape} benchmark.
Escape is blocked before $t_{\rm escape}$ and necessary afterward. The full
model best suppresses early false activation while still succeeding after
the escape opens. Main rows use $100$ paired episodes. $^\dagger$Low CVaR
with success ${<}0.21$ indicates low exposure from getting stuck; CIs are
in Appendix Table~\ref{tab:app_delayed_escape_ci}.}
\label{tab:delayed_escape}
\small
\setlength{\tabcolsep}{4pt}
\renewcommand{\arraystretch}{1.1}
\begin{tabular}{lcccc}
\toprule
Method & False pre-act$\downarrow$ & Suppress$\uparrow$
       & Success$\uparrow$ & Viol.\ CVaR$\downarrow$ \\
\midrule
Geometry-only policy  & $0.000$ & $1.000$ & $0.030$ & $1.894$ \\
Risk-loss-only          & $0.420$ & $0.580$ & $0.040$ & $1.793$ \\
Fixed-coeff context field     & $0.990$ & $0.010$ & $0.030$ & $0.463^\dagger$ \\
Black-box CVaR policy   & $0.920$ & $0.080$ & $0.200$ & $0.503^\dagger$ \\
DWA semantic            & $0.950$ & $0.050$ & $0.480$ & $0.695$ \\
MPPI semantic           & $0.950$ & $0.050$ & $0.240$ & $0.831$ \\
Ctx-enriched, expected cost  & $0.370$ & $0.630$ & $0.620$ & $0.855$ \\
\textbf{Route-aware Ctx CVaR}
  & $\mathbf{0.180}$ & $\mathbf{0.820}$
  & $\mathbf{0.810}$ & $\mathbf{0.740}$ \\
\bottomrule
\end{tabular}
\end{table}

\begin{figure*}[t]
\centering
\includegraphics[width=\textwidth]{NeurIPS_2026/figures/rellis_dyn_spotlight_four_methods.png}
\caption{\textbf{Qualitative temporal selectivity in one delayed-required
escape episode.} Yellow dashed trajectories/arrows show behavior before the
escape is available; solid colored trajectories/arrows show behavior after it
opens. The geometry-only policy ignores the material update, DWA and black-box CVaR move
before the escape is feasible and then stall, while route-aware context
enrichment suppresses before $t_{\rm escape}$ and takes the newly feasible
route after it opens.}
\label{fig:rellis_dyn_spotlight}
\end{figure*}

\emph{Hamiltonian structure}: black-box CVaR uses the same CVaR objective
and BEV patch but no Hamiltonian force channel; false pre-activation
$0.920$, success $0.200$ vs.\ $0.810$. Static CAR $0.884$ does not
transfer to temporal suppression.
\emph{CVaR training}: same architecture, different objective; CVaR halves
false pre-activation ($0.370{\to}0.180$), raises suppression
($0.630{\to}0.820$), improves success by $30\%$, reduces violation CVaR
($0.855{\to}0.740$).
\emph{Context-conditioning}: fixed vs.\ CNN-predicted $(\lambda_s,\lambda_h)$;
success $0.030$ vs.\ $0.810$. Without adaptive coefficients the method
suppresses nothing.

% ----------------------------------------------------------------------
\subsection{RELLIS-3D: Spatial Selectivity}
\label{sec:exp_rellis_static}
% ----------------------------------------------------------------------

Table~\ref{tab:rellis_selectivity} reports six variants on 2{,}250 BEV
episodes (leave-one-sequence-out; CIs in Appendix
Table~\ref{tab:app_rellis_full}). The SR gap between black-box CVaR ($1.292$) and route-aware context-enriched field
($2.358$) is the spatial analog of the temporal gap in
Table~\ref{tab:delayed_escape}: the black-box learns \emph{when} risk
gradients exist but not \emph{whether} following them is safe.
The route-aware gate produces the selectivity; Section~\ref{sec:exp_primary_delayed}
shows the same architecture without the gate (black-box) collapses to
$0.200$ navigation success on the strictest temporal test.

\begin{table}[t]
\centering
\caption{\textbf{RELLIS-3D spatial selectivity.}
Route-aware context enrichment achieves the best SR and AUPRC, indicating
selective activation rather than uniform risk-gradient following. Fixed
coefficients fail to suppress in blocked scenes (FAR~$0.888$).}
\label{tab:rellis_selectivity}
\small
\setlength{\tabcolsep}{5pt}
\renewcommand{\arraystretch}{1.1}
\begin{tabular}{lcccc}
\toprule
Method & CAR$\uparrow$ & FAR$\downarrow$ & SR$\uparrow$ & AUPRC$\uparrow$ \\
\midrule
Geometry-only policy    & $0.000$ & $0.000$ & $0.000$ & $0.000$ \\
Scalar context field            & $0.378$ & $0.752$ & $1.031$ & $0.008$ \\
Fixed-coeff context field    & $0.566$ & $0.888$ & $1.123$ & $0.008$ \\
Non-route directional Ctx  & $0.135$ & $0.178$ & $1.165$ & $0.206$ \\
Black-box CVaR policy     & $0.884$ & $0.129$ & $1.292$ & $0.206$ \\
\textbf{Route-aware Ctx} & $\mathbf{0.837}$ & $\mathbf{0.114}$
                           & $\mathbf{2.358}$ & $\mathbf{0.289}$ \\
\bottomrule
\end{tabular}
\end{table} 

% ----------------------------------------------------------------------
\subsection{RELLIS-Dyn: Temporal Selectivity}
\label{sec:exp_rellis_dyn}
% ----------------------------------------------------------------------

\paragraph{Material-event aggregate.}
Table~\ref{tab:rellis_dyn_agg} reports the three-event material subset
(mud onset, corridor closes, corridor opens) at 100 episodes per event.
Black-box CVaR success $0.273$ on the same task where it achieves
CAR~$0.884$ confirms that static selectivity metrics do not predict
navigation success.
Risk-loss-only hard exposure $0.694$ vs.\ Route-aware Ctx $0.027$
supports RQ6: the risk signal must enter the force channel.
The route-aware context-enriched field is the strongest zero-replan learned-field method,
matching planning baselines that use $9$--$68$ explicit replans.
Corridor-opens trajectory detail is in Appendix
Figure~\ref{fig:app_rellis_dyn_corridor}; the eight-event breakdown
(including moving-obstacle events where DWA is architecturally
stronger) is in Appendix Table~\ref{tab:app_rellis_dyn_8event_groups}.

\begin{table}[t]
\centering
\caption{\textbf{RELLIS-Dyn 3-event material aggregate.}
All methods receive the same updated BEV patch. The gate uses only
$K{=}8$ local primitives, not graph search or trajectory optimization.
$^\dagger$Oracle has global map access. $^{\ddagger}$Delay includes
gate-suppressed R2 episodes; see Table~\ref{tab:delayed_escape}.}
\label{tab:rellis_dyn_agg}
\small
\setlength{\tabcolsep}{3pt}
\renewcommand{\arraystretch}{1.1}
\resizebox{0.9\linewidth}{!}{%
\begin{tabular}{llccccccc}
\toprule
Method & Family
  & Succ.$\uparrow$ & Soft CVaR$\downarrow$
  & Delay$\downarrow$ & Stuck$\downarrow$
  & Replans$\downarrow$ & ms/step$\downarrow$ \\
\midrule
Geometry-only policy      & geom.-only     & $1.000$ & $0.712$ & $16.2$ & $0.04$ & $0$    & $0.8$  \\
Risk-loss-only     & learned loss   & $1.000$ & $0.673$ & $13.4$ & $0.03$ & $0$    & $0.9$  \\
Fixed-coeff Ctx     & learned field  & $0.120$ & $0.459$ & $1.47$ & $0.74$ & $0$    & $1.4$  \\
Black-box CVaR     & learned field  & $0.273$ & $0.482$ & $3.61$ & $0.61$ & $0$    & $4.2$  \\
Ctx-enriched, exp.cost  & learned field  & $0.983$ & $0.465$ & $1.30$ & $0.02$ & $0$    & $2.1$  \\
DWA semantic       & reactive       & $0.567$ & $0.621$ & $2.8$  & $0.43$ & $0$    & $0.6$  \\
CBF-QP             & safety filter  & $0.890$ & $0.651$ & $7.4$  & $0.09$ & $0$    & $0.9$  \\
Local A* (budget)  & planner        & $0.983$ & $0.629$ & $4.6$  & $0.01$ & $9.8$  & $8.4$  \\
MPC (budget)       & planner        & $0.967$ & $0.624$ & $3.8$  & $0.03$ & $11.1$ & $42.2$ \\
\textbf{Route-aware Ctx}& learned field & $\mathbf{0.980}$ & $\mathbf{0.638}$
  & $\mathbf{1.8}^{\ddagger}$ & $\mathbf{0.06}$ & $\mathbf{0}$ & $\mathbf{2.3}$ \\
Oracle replanner   & planner$^\dagger$  & $1.000$ & $0.601$ & $1.9$  & $0.00$ & $68.3$ & $97.5$ \\
\bottomrule
\end{tabular}}
\end{table}

\paragraph{Within-episode temporal selectivity.}
Figure~\ref{fig:rellis_dyn_spotlight} shows the delayed-required-escape
mechanism qualitatively.
Route-aware CVaR suppresses lateral activation throughout the blocked
window and activates within two steps of $t_{\rm escape}$
(false pre-activation $0.180$).
The expected-cost variant activates earlier ($0.370$), directly
demonstrating CVaR's tail-gradient contribution.
DWA false pre-activates on $95\%$ of episodes, typically $3$--$5$ steps
before $t_{\rm escape}$, because partial boundary changes trigger
its reactive response before the escape is fully open.

% ----------------------------------------------------------------------
\subsection{DFC2018: Geometry-Only Repair}
\label{sec:exp_dfc}
% ----------------------------------------------------------------------

Table~\ref{tab:dfc_geom_ctx} shows enrichment repairs a geometry-only
policy on a static semantic material map: success $0.867{\to}1.000$,
catastrophic failure $0.600{\to}0.100$, hard-hazard traversal $-96.9\%$,
cumulative risk $-59.2\%$, oscillation $-90.7\%$, path-length ratio
unchanged.
Full planner comparison in Appendix~\ref{app:dfc_full}; full-map
planners with global access obtain lower raw risk but the context-enriched field produces
smoother trajectories with zero oscillation.

\begin{table}[t]
\centering
\caption{\textbf{DFC2018 outcome shift} (300 paired episodes).}
\label{tab:dfc_geom_ctx}
\scriptsize\setlength{\tabcolsep}{3pt}\renewcommand{\arraystretch}{1.08}
\resizebox{0.7\linewidth}{!}{%
\begin{tabular}{lrrrrrr}
\toprule
Method & Succ.$\uparrow$ & Cat.fail$\downarrow$ & Hard len.$\downarrow$
       & Risk$\downarrow$ & Len.ratio & Osc.$\downarrow$ \\
\midrule
Geometry-only policy & $0.867$ & $0.600$ & $21.106$ & $23.267$ & $1.044$ & $82.212$ \\
\textbf{Ctx-enriched} & $\mathbf{1.000}$ & $\mathbf{0.100}$ & $\mathbf{0.658}$
                 & $\mathbf{9.493}$ & $1.045$ & $\mathbf{7.629}$ \\
\bottomrule
\end{tabular}}
\end{table}

% ----------------------------------------------------------------------
\subsection{Highway-env: Interaction Selectivity}
\label{sec:exp_highway}
% ----------------------------------------------------------------------

The open slow-leader scenario is the interaction analog of corridor
opens (lateral channel should activate); the boxed scenario is the
analog of moving-obstacle-blocks-detour (channel should suppress).
Table~\ref{tab:highway_mechanism} shows the mechanism ablation: removing
the lateral channel causes collision with the slow leader on every
episode; adding it without the TTC channel causes off-road failure in
the boxed scenario; only the full system passes and suppresses correctly.
MOBIL-IDM and risk-aware MPC outperform the context-enriched field overall (Appendix
Table~\ref{tab:app_highway_baselines}); the context-enriched field is the strongest
learned method and the only one exhibiting the intended
activation/suppression split.

\begin{table}[t]
\centering
\caption{\textbf{Highway mechanism ablation} (200 paired seeds).
The lateral channel enables passing; the TTC channel prevents using it
when the adjacent maneuver is unsafe.}
\label{tab:highway_mechanism}
\scriptsize\setlength{\tabcolsep}{3pt}\renewcommand{\arraystretch}{1.08}
\resizebox{\linewidth}{!}{%
\begin{tabular}{lrrrrrr}
\toprule
Variant & Def.off-rd$\downarrow$ & Slow coll$\downarrow$
        & Slow succ$\uparrow$ & Boxed coll$\downarrow$
        & Boxed off-rd$\downarrow$ & Boxed spd$\uparrow$ \\
\midrule
No lateral context       & $0.00$ & $1.00$ & $0.00$ & $1.00$ & $0.00$ & $22.33$ \\
Lateral only             & $0.20$ & $0.00$ & $1.00$ & $1.00$ & $0.00$ & $22.33$ \\
\textbf{Lateral + TTC}   & $\mathbf{0.00}$ & $\mathbf{0.00}$ & $\mathbf{1.00}$
                         & $\mathbf{0.00}$ & $\mathbf{0.00}$ & $\mathbf{7.57}$ \\
\bottomrule
\end{tabular}}
\end{table}

% ----------------------------------------------------------------------
\subsection{Interpretation and Limitations}
\label{sec:exp_interpretation}
% ----------------------------------------------------------------------

\paragraph{Tradeoff positions.}
Reactive baselines (DWA, CBF-QP) handle hard-boundary events well but
false pre-activate on discovery events (DWA: $0.950$ on delayed escape).
Planning baselines reduce violation under replan budget but pay
$9$--$68\times$ more replans.
The route-aware context-enriched field occupies the zero-replan gradient-following frontier:
strongest on soft-risk and escape-discovery events, weaker on
moving-obstacle events where reactive boundary following suffices.
The eight-event Pareto (Appendix Figure~\ref{fig:app_rellis_dyn_pareto})
shows the full tradeoff surface.

\paragraph{Limitations.}
(i) $b_{\rm sp}(\phi)$ is a differentiable barrier, not a CBF
certificate; our claims are empirical tail-risk reductions, not
forward-invariance guarantees
(Appendix~\ref{app:theory}, \S A.9).
(ii) Performance depends on risk-map quality; at $20$--$30\%$ label
corruption, CAR/FAR degrade despite the gate's tendency to suppress
inconsistent soft-risk evidence (Appendix
Table~\ref{tab:app_rellis_perception}).
(iii) Gate overhead must be profiled on latency-constrained hardware
(ms/step in Table~\ref{tab:rellis_dyn_agg}).


% \clearpage
